% TOP_PROBLEM: {"id": 8900008, "title": "Congruent number decision problem", "category_id": 3, "category": {"id": 3, "name": "graph_theory", "display_name": "Graph Theory", "description": "Problems involving graphs, networks, and their properties.", "slug": "graph-theory", "order_index": 3, "created_at": "2026-07-31T15:26:25.671Z"}, "status": "open"}
% FIRST_POSTED: 2026-09-25
% ENTRY_KIND: statement_only
% !TeX program = lualatex
% Definition and sources only; this file is not an LLM solution attempt.
\documentclass[11pt]{article}
\usepackage[margin=25mm]{geometry}
\usepackage{fontspec}
\setmainfont{Latin Modern Roman}
\usepackage{amsmath,amssymb,mathtools}
\usepackage{unicode-math}
\setmathfont{Latin Modern Math}
\usepackage{xurl,hyperref}
\hypersetup{colorlinks=true,urlcolor=blue}
\setcounter{secnumdepth}{0}
\setlength{\parindent}{0pt}
\setlength{\parskip}{5pt}
\setlength{\emergencystretch}{3em}
\begin{document}
\section*{\texorpdfstring{Congruent number decision problem}{Congruent number decision problem}}\label{8900008}
\subsection{Definitions and mathematical statement}
A positive integer $n$ is congruent if there exist $a,b,c\in{\mathbb{Q}}_{>0}$ with
\[
 a^2+b^2=c^2,\qquad ab/2=n.
\]
Equivalently, the elliptic curve $E_n:y^2=x^3-n^2x$ has a rational point with $y\ne0$. For such a point one obtains side lengths $|x^2-n^2|/|y|$, $2n|x|/|y|$, and $(x^2+n^2)/|y|$. Multiplication of $n$ by a rational square does not change congruence, so positive integers reduce to squarefree parts. The problem seeks an unconditional exact characterization, not only a conditional criterion.
\par\smallskip\textit{Formulation and concept references:} \href{https://kconrad.math.uconn.edu/blurbs/ugradnumthy/congnumber.pdf}{[S1]}.

\subsection{Short English statement}
Which positive integers occur as the areas of right triangles with entirely rational side lengths?
\subsection{Sources}
\begin{itemize}
\item[S1]Keith Conrad, "The Congruent Number Problem," expository notes, University of Connecticut, kconrad.math.uconn.edu/blurbs/ugradnumthy/congnumber.pdf.
\url{https://kconrad.math.uconn.edu/blurbs/ugradnumthy/congnumber.pdf}.
\textit{Formulation source.}
\end{itemize}
\end{document}
